\color{uniblue}\subsection[Выбор уставок максимальной токовой защиты]{Выбор уставок максимальной токовой защиты}
\color{black}

\begin{enumerate}
\item
Первичный ток срабатывания ступени $I_{\textsl{с.з.МТЗ}}^I$ определяется по условию отстройки от максимального рабочего тока трансформатора по выражению (\ref{eq:mtz1}):
\begin{equation*}
    I_{\textsl{с.з.МТЗ}}^I = \frac {k_{\textsl{н}} \cdot k_{\textsl{перег}}}{k_{\textsl{в}}}I_{\textsl{раб.макс.тр}} = \frac{1,2 \cdot 2,0}{0,95} \cdot 1,05 \cdot 250 = 663\; (\textsl{А}) ,
\end{equation*}

\item
Выдержка времени выбирается на ступень селективности больше времени срабатывания самой медленно действующей защиты питающей сети по выражению (\ref{eq:mtz4}):
\begin{equation*}
t_{\textsl{с.з.МТЗ}}^I = t_{\textsl{с.з.МТЗ}}^{\textsl{фид}}+{\Delta t} = 1,6 + 0,3 = 1,9\; (\textsl{с}) ,
\end{equation*}

\item
Чувствительность ступени МТЗ определяется по току двухфазного КЗ в минимальном режиме работы сети на шинах НН по выражению (\ref{eq:mtz5}):
\begin{equation*}
k_{\textsl{ч}} = \frac{I_{\textsl{к.мин.}}^{(2)}}{I_{\textsl{с.з. МТЗ}}^{I}} = \frac{0,87 \cdot 840}{663} = 1,1 < 1,3,
\end{equation*}
\item
Ступень МТЗ не удовлетворяет не требованиям по чувствительности, при недостаточной чувствительности используется пуск по напряжению со стороны НН трансформатора, расчет тока срабатывания ступени МТЗ в этом случае происходит по выражению (\ref{eq:mtz6}).
\newpage
\begin{equation*}
I_{\textsl{с.з.МТЗ}}^{I} = \frac {k_{\textsl{н}}}{k_{\textsl{в}}}I_{\textsl{раб.макс.тр}} = \frac{1,2}{0,95} \cdot 1,05 \cdot 250 = 332\; (\textsl{А}),
\end{equation*}
\item
Чувствительность ступени МТЗ в этом случае определяется по выражению (\ref{eq:mtz5}):
\begin{equation*}
k_{\textsl{ч}} = \frac{I_{\textsl{к.мин.}}^{(2)}}{I_{\textsl{с.з. МТЗ}}^{I}} = \frac{0,87 \cdot 840}{332} = 2,2 > 1,3,
\end{equation*}
Ступень МТЗ удовлетворяет требованиям по чувствительности.

Уставка задается в относительных единицах от номинального тока аналогового входа устройства, для этого полученное значение пересчитывается по следующему выражению:
\begin{equation*}
    I_{\textsl{ср}} = \frac{I_{\textsl{с.з.МТЗ}}^I}{I_{\textsl{ном.вх.уст}}} \cdot \frac{I_{\textsl{ном.вт.ТТ}}}{I_{\textsl{ном.пер.ТТ}}} = \frac{332}{5} \cdot \frac{5}{400} = 0,83\; (\textsl{о.е.}),
\end{equation*}
\begin{eqexpl}[15mm]
    \item{$I_{\textsl{ном.вх.уст}}$} номинальный ток аналоговых входов устройства, 1 или 5 А;
    \item{$I_{\textsl{ном.вт.ТТ}}$} номинальный вторичный ток ТТ, А;
    \item{$I_{\textsl{ном.пер.ТТ}}$} номинальный первичный ток ТТ, А.
\end{eqexpl}
\item
Параметр \textbf{$I_{\textsl{ср}}$} <<Ток срабатывания>> ступени МТЗ принимается \textbf{равным 0,83}.
\item
Параметр \textbf{$T_{\textsl{ср}}$} <<Выдержка времени срабатывания>> ступени МТЗ принимается \textbf{равным 1,9 c}.

\item 
Минимальное рабочее напряжение на стороне НН трансформатора определяется по выражению (\ref{eq:mtz7}):
\begin{equation*}
U_{\textsl{с.з.НН}} = \frac {U_{\textsl{р.мин.}}}{k_{\textsl{н}} \cdot k_{\textsl{в}} \cdot k_{\textsl{сам}}} =\frac{0,6 \cdot 6000}{1,2 \cdot 1,05 \cdot 1,0} = 2857\; (\textsl{В}),
\end{equation*}
Для ввода в устройство необходимо пересчитать значение напряжения во вторичные величины по выражению:
\begin{equation*}
    U_{\textsl{с.р.}} = U_{\textsl{с.з.НН}} \cdot \frac{U_{\textsl{ном.вт.ТН}}}{U_{\textsl{ном.пер.ТН}}} = 2857 \cdot \frac{100}{6000} = 47,6\; (\textsl{В}),
\end{equation*}
\item
Параметр \textbf{$U_{\textsl{ср}}$} <<Напряжение срабатывания>> функции КПОН принимается \textbf{равным 47,6 В}.

\end{enumerate}